\documentclass[a4paper,11pt]{article}
\usepackage{iftex}
\ifPDFTeX
\usepackage[T1]{fontenc}
\usepackage[utf8]{inputenc}
\usepackage{lmodern}
\else
\usepackage{fontspec}
\setmainfont{Latin Modern Roman}
\fi
\usepackage[ngerman]{babel}
\usepackage{amsmath,amssymb,amsthm,mathtools}
\usepackage{booktabs,tabularx,array}
\usepackage[a4paper,top=24mm,bottom=25mm,left=27mm,right=27mm]{geometry}
\usepackage{microtype}
\usepackage[hidelinks]{hyperref}
\newcommand{\Pp}[1]{\mathcal P_+(#1)}
\newcommand{\Sing}{\operatorname{Sing}}
\newtheorem{satz}{Satz}
\newtheorem{lemma}[satz]{Hilfssatz}
\newtheorem{korollar}[satz]{Folgerung}
\setlength{\parindent}{0pt}
\setlength{\parskip}{.55em plus .1em minus .1em}
\setlength{\emergencystretch}{1.2em}
\allowdisplaybreaks[2]
\hypersetup{pdftitle={Band 37 - Kontextprofile und Induktion - Lesefassung},
 pdfauthor={Martin Kunze},pdfsubject={Bewiesene Rekonstruktionssaetze und Grenzen des allgemeinen Induktionsschritts}}
\title{Kontextprofile und die Potenzhalbgruppenvermutung\\[.5em]
\large Induktionsanfang, Rekonstruktionssätze und Grenzen\\Lesefassung zu Band 37}
\author{Martin Kunze}
\date{6. Oktober 2026}
\begin{document}
\maketitle

Aus ausgewählten Einermengen von Produktwerten lassen sich ein
Kongruenzquotient der Grundhalbgruppe und die Größe jeder seiner
Klassen zurückgewinnen. Trennt dieser Quotient die Produktwerte,
rekonstruiert er bereits die ganze Halbgruppe. Dieser
Rekonstruktionssatz wird im Folgenden vollständig bewiesen.

Für die allgemeine Induktion bleibt jedoch eine wesentliche Lücke:
Die trennenden Markierungen müssen erst verfügbar sein. Für einen
beliebig vorgegebenen Potenzisomorphismus ist das sogar falsch.
Auch ein kleinerer Quotient mit bekannten Klassengrößen erlaubt
allein noch keine Rückgewinnung der ursprünglichen Multiplikation.
Wir beweisen beide Grenzen durch konkrete endliche Beispiele.

Der bewiesene Zweiproduktfall dient als Induktionsanfang. Der neue
Profilsatz liefert außerdem einen kurzen vollständigen Beweis des
Dreiproduktfalls mit verschwindenden Dreifachprodukten. Ein allgemeiner
Induktionsbeweis für sämtliche endlichen Halbgruppen folgt daraus
bislang nicht.

Wir verwenden ausschließlich nichtleere Teilmengen:
\[
 \Pp C=\{A\subseteq C:A\ne\varnothing\},\qquad
 AB=\{ab:a\in A,\ b\in B\}.
\]
Alle Grundhalbgruppen dieser Lesefassung sind endlich und nichtleer.

\section{Der bewiesene Induktionsanfang}

Die gewünschte Aussage für eine natürliche Zahl \(k\ge1\) lautet
\[
 \mathsf H(k):\quad
 \forall C,D\quad
 \bigl(|C^2|\le k\ \land\ \Pp C\cong\Pp D\bigr)
 \ \Longrightarrow\ C\cong D.
\]
Hier durchlaufen \(C,D\) die endlichen nichtleeren Halbgruppen.
Ein zunächst beliebiges \(D\) wäre unter der Isomorphieannahme
ohnehin endlich, weil \(d\mapsto\{d\}\) es in \(\Pp D\) einbettet.

Für \(|C^2|=1\) ist die Multiplikation konstant. Das konstante
Produktelement ist durch die Multiplikation ausgezeichnet; die
restlichen Elemente können beliebig bijektiv zugeordnet werden.
Ferner ist \(|\Pp C|=2^{|C|}-1\), sodass ein Potenzisomorphismus
die Grundgröße erhält. Auch das Produktquadrat von \(D\) muss
einelementig sein: Seine Produkt-Einermengen liegen im einelementigen
Quadrat von \(\Pp D\). Damit ist \(\mathsf H(1)\) bewiesen.

Für \(|C^2|=2\) verwenden wir den bereits bewiesenen Satz aus der
Lesefassung \emph{Zwei Produktwerte genügen}, Band~37, vom
5.~Oktober 2026. Dort wird ohne Voraussetzung an \(C^3\) gezeigt:
\[
 |C^2|=2,\quad \Pp C\cong\Pp D
 \quad\Longrightarrow\quad C\cong D.
\]
Zusammen mit dem konstanten Fall steht damit
\(\mathsf H(2)\) als Induktionsanfang fest.

Für \(k\ge3\) müsste aus \(\mathsf H(k-1)\) die Aussage
\(\mathsf H(k)\) folgen. Dazu genügt es nicht, nur irgendeinen
kleineren Quotienten zu konstruieren: Es muss auch bewiesen werden,
wie dessen Isomorphismus auf die ursprünglichen Halbgruppen gehoben
werden kann.

\section{Kontexte unterscheiden ausgewählte Produktwerte}

Sei \(E\subseteq C^2\) eine Menge markierter Produktwerte.
Wir schreiben \(C^\sharp=C\mathbin{\dot\cup}\{\mathbf1\}\) für die
Halbgruppe mit einer \emph{zusätzlichen formalen Eins}, auch wenn \(C\)
bereits eine Eins besitzt. Der Faktor \(\mathbf1\) bedeutet, dass auf
dieser Seite kein Faktor steht. Definiere
\begin{equation}\label{cp:point}
 a\equiv_E b
 \quad\Longleftrightarrow\quad
 \forall x,y\in C^\sharp\ \forall e\in E:
 \quad (xay=e\ \Longleftrightarrow\ xby=e).
\end{equation}
Insbesondere ist der vollständig leere Kontext \(x=y=\mathbf1\)
zugelassen.

Diese Äquivalenzrelation ist eine Kongruenz. Ist etwa \(a\equiv_E b\)
und \(z\in C\), dann ersetzt der Vergleich von \(za\) und \(zb\)
den linken Kontext \(x\) nur durch \(xz\). Für \(az,bz\) gilt das
entsprechende Argument rechts. Also ist der Quotient
\(C/{\equiv_E}\) eine Halbgruppe.

Nun fügen wir auch \(\Pp C\) eine zusätzliche formale Eins hinzu.
Für \(A\in\Pp C\) sei
\[
 \tau_C(A)=
 \bigl\{(U,V,e):
 U,V\in(\Pp C)^\sharp,\ e\in E,\ UAV=\{e\}\bigr\}.
\]
Das Profil \(\tau_C(A)\) verwendet nur die Multiplikation und die
markierten Elemente \(\{e\}\). Die Mengendarstellung der übrigen
Elemente muss dafür nicht bekannt sein.

Weil alle echten Faktoren nichtleer sind, gilt
\[
 UAV=\{e\}
 \quad\Longleftrightarrow\quad
 \forall a\in A:\ U\{a\}V=\{e\}.
\]
Daher erhalten wir die grundlegende Identität
\begin{equation}\label{cp:intersection}
 \tau_C(A)=\bigcap_{a\in A}\tau_C(\{a\}).
\end{equation}
Leere äußere Kontexte ändern an dieser Aussage nichts.

\section{Der Quotient und alle Klassengrößen sind rekonstruierbar}

\begin{satz}[Profilrekonstruktion]\label{cp:profiles}
Seien \(\Phi:\Pp C\to\Pp D\) ein Isomorphismus,
\(E\subseteq C^2\), \(F\subseteq D^2\) und \(h:E\to F\)
eine Bijektion mit
\[
 \Phi(\{e\})=\{h(e)\}\qquad(e\in E).
\]
Dann induziert \(\Phi\) einen Isomorphismus
\[
 \theta:C/{\equiv_E}\longrightarrow D/{\equiv_F}
\]
der Kongruenzquotienten. Entsprechende Klassen besitzen dieselbe
Anzahl ursprünglicher Elemente.
\end{satz}

\begin{proof}
Betrachte die endliche Menge aller auftretenden Profile, geordnet
durch Inklusion. Für ein solches Profil \(p\) setzen wir
\[
 M(p)=|\{A\in\Pp C:\tau_C(A)\supseteq p\}|,\qquad
 D_p=\{a\in C:\tau_C(\{a\})\supseteq p\}.
\]
Nach \eqref{cp:intersection} sind die bei \(M(p)\) gezählten Mengen
genau die nichtleeren Teilmengen von \(D_p\). Also gilt
\begin{equation}\label{cp:count}
 M(p)=2^{N(p)}-1,\qquad N(p)=|D_p|=\log_2(M(p)+1).
\end{equation}
Der Logarithmus bezeichnet hier den eindeutig bestimmten endlichen
Exponenten.

Sei \(m(p)\) die Anzahl der Grundelemente mit genau dem Profil \(p\).
Dann gilt
\begin{equation}\label{cp:mobius}
 m(p)=N(p)-\sum_{q\supsetneq p}m(q).
\end{equation}
Die Summe läuft über die tatsächlich auftretenden Profile. Beginnend
mit den maximalen Profilen bestimmt diese absteigende Rekursion alle
Zahlen \(m(p)\). Genau die Profile mit \(m(p)>0\) sind Profile
ursprünglicher Einzelpunkte. Andere auftretende Profile dürfen die
Zahl null erhalten und werden anschließend weggelassen.

Auch die Multiplikation der Profile ist erkennbar. Profilgleichheit
ist eine Kongruenz auf \(\Pp C\): Aus
\(\tau_C(A)=\tau_C(B)\) folgt etwa
\(\tau_C(AZ)=\tau_C(BZ)\), weil
\[
 U(AZ)V=UA(ZV).
\]
Der zusätzliche Faktor wird lediglich in den rechten Kontext
aufgenommen. Links verfährt man genauso. Folglich ist
\[
 p\cdot q:=\tau_C(AB)
 \quad\text{für }\tau_C(A)=p,\ \tau_C(B)=q
\]
unabhängig von der Wahl der Repräsentanten.

Die Profile mit \(m(p)>0\) bilden eine Unterhalbgruppe dieses
Profilquotienten, denn
\[
 \tau_C(\{a\})\cdot\tau_C(\{b\})=\tau_C(\{ab\}).
\]
Zwei Einzelpunktprofile sind genau dann gleich, wenn
\(a\equiv_E b\) gilt. Eine Richtung erhält man durch Einermengen
als äußere Kontexte. Für die andere Richtung ist eine Gleichung mit
Mengenkontexten die gemeinsame Erfüllung aller entsprechenden
Einzelkontextgleichungen. Damit ist die Unterhalbgruppe der
Einzelpunktprofile genau \(C/{\equiv_E}\), und ihre Gewichte \(m(p)\)
sind die Klassengrößen.

Schließlich erweitern wir \(\Phi\) auf die formale Eins durch
\(\Phi^\sharp(\mathbf1)=\mathbf1\). Die Kontextumbenennung
\[
 (U,V,e)\longmapsto
 \bigl(\Phi^\sharp(U),\Phi^\sharp(V),h(e)\bigr)
\]
überführt \(\tau_C(A)\) in \(\tau_D(\Phi(A))\). Sie erhält
Profilinklusion, Profilmultiplikation und die Zahlen \(M(p)\).
Deshalb erhält sie nach \eqref{cp:count} und \eqref{cp:mobius}
auch \(N(p)\) und \(m(p)\). Die erkannten Einzelpunktprofile werden
somit samt Multiplikation und Gewichten bijektiv zugeordnet.
\end{proof}

Dieser Satz liefert den Quotientenisomorphismus unmittelbar aus
\(\Phi\). Für die Rekonstruktion des Quotienten selbst wird keine
Induktionsannahme benötigt.

\section{Wann die rekonstruierte Information genügt}

\begin{satz}[Trennung der Produktwerte]\label{cp:separation}
Unter den Voraussetzungen von Satz~\ref{cp:profiles} enthalte jede
\(\equiv_E\)-Klasse höchstens ein Element aus \(C^2\), und jede
\(\equiv_F\)-Klasse höchstens ein Element aus \(D^2\). Dann ist
\(C\cong D\).
\end{satz}

\begin{proof}
Seien \(\pi_C,\pi_D\) die beiden Quotientenabbildungen und
\(\theta\) der Isomorphismus aus Satz~\ref{cp:profiles}. Es gilt
\[
 (C/{\equiv_E})^2=\pi_C(C^2).
\]
Eine Klasse gehört also genau dann zum Quadrat des Quotienten,
wenn sie ein ursprüngliches Produktelement enthält. Entsprechendes
gilt auf der Zielseite.

Da \(\theta\) die Quadrate der Quotienten zuordnet, besitzen
entsprechende Produktklassen nach Voraussetzung jeweils genau ein
ursprüngliches Produktelement. Wähle zwischen jeder Paarung gleich
großer Klassen eine Bijektion, die diese besonderen Elemente
aufeinander abbildet. Die restlichen Elemente können innerhalb der
Klassen beliebig bijektiv zugeordnet werden. Die Vereinigung ist eine
Bijektion \(f:C\to D\) mit
\[
 \pi_D(f(a))=\theta(\pi_C(a)).
\]
Für \(a,b\in C\) folgt
\[
 \pi_D(f(ab))
 =\theta(\pi_C(a)\pi_C(b))
 =\pi_D(f(a)f(b)).
\]
Sowohl \(f(ab)\) als auch \(f(a)f(b)\) ist ein Produktelement von
\(D\). Das erste ist es durch die vorgeschriebene Wahl von \(f\),
das zweite durch seine Form. In ihrer gemeinsamen Klasse gibt es
genau ein solches Element. Daher sind sie gleich, und \(f\) ist
multiplikativ.
\end{proof}

\begin{korollar}\label{cp:kminusone}
Gilt \(|C^2|=|D^2|=k\ge2\), so genügt es, dass ein
Potenzisomorphismus \(k-1\) verschiedene Produkt-Einermengen von
\(C\) auf Produkt-Einermengen von \(D\) abbildet.
\end{korollar}

\begin{proof}
Markiere diese \(k-1\) Produktwerte. Jeder markierte Wert \(e\)
bildet eine einelementige Kontextklasse: Der leere Kontext prüft
bereits die Gleichung \(a=e\). Unter den übrigen Elementen befindet
sich nur noch ein einziger Produktwert. Also enthält auch jede
unmarkierte Klasse höchstens einen Produktwert. Dasselbe gilt auf
der Zielseite. Satz~\ref{cp:separation} ist anwendbar.
\end{proof}

Für \(k=2\) ergibt dies zusätzlich eine kurze zweite Begründung
unseres Induktionsanfangs: Das Quadrat von \(\Pp C\) besteht aus
\(\{a\},\{b\},C^2\). Daraus folgt auch \(|D^2|=2\). Bei nur einem
Produktwert hätte das Zielquadrat ein Element; bei mindestens drei
Produktwerten hätte es mindestens drei Einermengen und zusätzlich
\(D^2\), also mindestens vier Elemente.

Von den zwei Bildern der Produkt-Einermengen muss mindestens eines
eine Einermenge sein, weil im dreielementigen Zielquadrat nur ein
Element keine Einermenge ist. Folgerung~\ref{cp:kminusone} schließt
den Beweis. Der bereits veröffentlichte Zweiproduktsatz bleibt damit
ein eigenständig bewiesener Anfangsfall.

\section{Ein kurzer weiterer Anfangsfall}

Die Profilrekonstruktion vereinfacht auch den bereits bewiesenen
Fall \(C^2=\{0,u,v\}\), \(C^3=\{0\}\). Zunächst benötigen wir eine
Schranke, die ohne Nilpotenzannahme gilt.

\begin{lemma}\label{cp:square-bound}
Für \(k=|C^2|\ge3\) gilt
\[
 2k\le|(\Pp C)^2|\le2^k-1.
\]
Insbesondere erzwingt ein Potenzisomorphismus bei \(|C^2|=3\)
auch \(|D^2|=3\).
\end{lemma}

\begin{proof}
Die obere Schranke folgt daraus, dass alle Mengenprodukte nichtleere
Teilmengen von \(C^2\) sind. Für die untere Schranke betrachte einen
Graphen mit Knotenmenge \(C^2\). Verbinde zwei verschiedene Produktwerte,
wenn sie in derselben Zeile oder Spalte der Multiplikationstafel
vorkommen.

Jede Kante \(\{c,d\}\) ergibt ein zweielementiges Mengenprodukt:
Bei gemeinsamer Zeile wählt man einen linken Einzelpunkt und die
beiden rechten Faktoren, bei gemeinsamer Spalte umgekehrt.
Der Graph ist zusammenhängend. Denn zwischen \(c=xy\) und \(d=zt\)
vermittelt der Wert \(xt\): Der erste Übergang benutzt eine gemeinsame
Zeile, der zweite eine gemeinsame Spalte. Ist ein Übergang trivial,
kann er weggelassen werden.

Ein zusammenhängender Graph auf \(k\) Knoten hat mindestens \(k-1\)
Kanten. Also gibt es mindestens \(k-1\) verschiedene zweielementige
Mengenprodukte. Hinzu kommen die \(k\) Produkt-Einermengen und
\(CC=C^2\), das wegen \(k\ge3\) von ihnen allen verschieden ist.
Zusammen sind dies mindestens \(2k\) Mengenprodukte.

Für \(k=3\) besitzt das Potenzquadrat somit sechs oder sieben
Elemente. Ein Zielquadrat mit mindestens vier Grundproduktwerten
hätte nach der unteren Schranke mindestens acht Elemente.
Bei einem beziehungsweise zwei Grundproduktwerten wären es nur
eins beziehungsweise drei. Daher ist \(|D^2|=3\).
\end{proof}

\begin{satz}\label{cp:three}
Ist \(C\) endlich mit \(C^2=\{0,u,v\}\),
\(C^3=\{0\}\) und drei verschiedenen Produktwerten, so folgt aus
\(\Pp C\cong\Pp D\) die Isomorphie \(C\cong D\).
\end{satz}

\begin{proof}
Hilfssatz~\ref{cp:square-bound} liefert \(|D^2|=3\).
Alle Dreifachprodukte in \(\Pp C\) sind \(\{0\}\). Ihre Bilder
sind daher dasselbe Element \(Z=\Phi(\{0\})\). Testet man auf
der Zielseite drei Einermengen, so sieht man, dass \(Z=\{z\}\)
und \(D^3=\{z\}\) gilt. Jedes Vierfachprodukt ist ein Dreifachprodukt
nach Zusammenfassen zweier Faktoren, also ist \(D^4=\{z\}\).
Schreibt man \(z=abc\), folgt für jedes \(t\in D\)
\(tz=tabc=z\) und \(zt=abct=z\). Somit ist \(z=0_D\) die Null
von \(D\). Dasselbe Argument rechtfertigt die Null auf der Quellseite.
Insbesondere wird \(\{0_C\}\) auf \(\{0_D\}\) abgebildet.

Für ein echtes Einermengenziel \(\{c\}\) hat die Gleichung
\(AB=\{c\}\), bei festem \(A\), stets \(2^m-1\) Lösungen:
Die Lösungen sind die nichtleeren Teilmengen der Menge aller
\(b\) mit \(ab=c\) für jedes \(a\in A\).

Eine Produktmenge \(Z\), die \(0_D\) und mindestens einen
Nichtnullwert enthält, besitzt dagegen für einen geeigneten
linken Faktor eine positive gerade Lösungszahl. Schreibe \(Z=A_0B_0\)
und setze \(A=A_0\cup\{0_D\}\). Dann ist \(AB_0=Z\).
In jeder Lösung \(B\) darf man die Zugehörigkeit von \(0_D\)
umkehren, ohne das Produkt oder die Nichtleerheit zu ändern.
Denn ein Nichtnullprodukt verlangt einen rechten Faktor
verschieden von \(0_D\), und die Null im Produkt wird durch \(A\)
bereits erzeugt. Die Lösungen werden also durch eine fixpunktfreie
Involution gepaart.

Solche Produktmengen können folglich keine Bilder von
Produkt-Einermengen sein. Schreiben wir
\(D^2=\{0_D,u_D,v_D\}\), so liegen die beiden verschiedenen Bilder
von \(\{u\},\{v\}\) unter
\[
 \{u_D\},\quad\{v_D\},\quad\{u_D,v_D\}.
\]
Mindestens eines ist eine Einermenge. Zusammen mit der erhaltenen
Null-Einermenge sind somit zwei der drei Produkt-Einermengen
zugeordnet. Folgerung~\ref{cp:kminusone} liefert \(C\cong D\).
\end{proof}

Dieser Beweis kommt ohne die längere Begleiterzählung der bisherigen
Dreiprodukt-Lesefassung aus. Er betrifft weiterhin den genannten
Teilfall von \(\mathsf H(3)\), nicht schon ganz \(\mathsf H(3)\).

\section{Ein kleiner Quotient bestimmt seine Erweiterung nicht}

Das folgende Beispiel trifft genau die Stelle, an der die allgemeine
Induktion noch keine Folgerung liefert. Betrachte auf
\(\{0,u,v,a,b\}\) die beiden Halbgruppen mit den aktiven Tabellen
\[
 \begin{array}{c|cc}
 C_1&a&b\\\hline
 a&u&u\\
 b&v&v
 \end{array}
 \qquad\qquad
 \begin{array}{c|cc}
 C_2&a&b\\\hline
 a&u&v\\
 b&v&u
 \end{array}.
\]
Alle Produkte mit mindestens einem Faktor aus \(\{0,u,v\}\)
seien \(0\). Jede Klammerung eines Dreifachprodukts ergibt daher
Null; beide Tabellen sind assoziativ. Ferner gilt
\[
 C_i^2=\{0,u,v\},\qquad C_i^3=\{0\}.
\]

Markieren wir \(E=\{0\}\), entstehen in beiden Halbgruppen
genau die Klassen
\begin{equation}\label{cp:example-classes}
 O=\{0\},\qquad R=\{u,v\},\qquad X=\{a,b\}.
\end{equation}
Der leere Kontext trennt \(0\) von allen anderen Elementen.
Die Werte \(u,v\) sind nicht zu trennen: Jeder nichtleere äußere
Kontext macht sie zu Null. Bei \(a,b\) hat ein einseitiger Kontext
genau dann ein Nullprodukt, wenn der äußere Faktor in
\(\{0,u,v\}\) liegt. Jeder beidseitige echte Kontext ergibt wieder
Null. Das trennt \(X\) von \(R\), unterscheidet innerhalb von \(X\)
aber keine Punkte.

Die Gewichte sind in beiden Fällen \(1,2,2\). Der Quotient hat
die Multiplikation
\[
 Q=\{O,R,X\},\qquad X^2=R,\qquad
 \text{alle übrigen Produkte sind }O.
\]
Es gilt also \(|Q^2|=2\). Wir gelangen tatsächlich unmittelbar zum
bereits bewiesenen Induktionsanfang.

Sogar die Profile sämtlicher nichtleerer Teilmengen stimmen unter
der natürlichen Benennung der Kontexte überein. Ihre drei Werte
bilden die Kette
\[
 \tau(\{0\})\supsetneq\tau(\{u\})=\tau(\{v\})
 \supsetneq\tau(\{a\})=\tau(\{b\}).
\]
Wegen der Schnittformel sind die zugehörigen Anzahlen
\[
 1,\qquad 2^3-2=6,\qquad 2^5-2^3=24.
\]
Die erste Familie besteht nur aus \(\{0\}\), die zweite aus den
übrigen nichtleeren Teilmengen von \(\{0,u,v\}\), die dritte aus
allen Mengen, die \(a\) oder \(b\) enthalten. Auch ihre
Profilmultiplikation ist diejenige von \(Q\).

Trotzdem sind \(C_1,C_2\) nicht isomorph: In \(C_1\) gilt
\(ab=u\ne v=ba\), während \(C_2\) kommutativ ist.
Der gewichtete Quotient und sogar der vollständige Profilquotient
mit allen Profilfasergrößen bestimmen die Erweiterung somit nicht.

Dies ist kein Gegenbeispiel zur Potenzhalbgruppenvermutung.
Auch \(\Pp{C_1}\) ist nicht kommutativ und \(\Pp{C_2}\) ist
kommutativ, sodass ihre Potenzhalbgruppen gerade nicht isomorph
sind. Das Beispiel widerlegt die Rekonstruktion aus den bereits
ausgewerteten Quotientendaten allein. Ein allgemeiner
Erweiterungssatz müsste weitere Information des ursprünglichen
\(\Phi\) benutzen.

\section{Die Trennung ist für ein gegebenes Bild nicht garantiert}

Sei \(L=\{a,b,c\}\) eine Linksnullhalbgruppe, also \(xy=x\).
Dann gilt \(L^2=L\) und für alle nichtleeren Teilmengen
\[
 AB=A.
\]
Jede Permutation der sieben Elemente von \(\Pp L\) ist daher
ein Automorphismus. Insbesondere gilt dies für die Involution
\[
 \Phi(A)=
 \begin{cases}
 L\setminus A,&\varnothing\ne A\subsetneq L,\\
 L,&A=L.
 \end{cases}
\]
Sie vertauscht alle drei Einermengen mit Zweiermengen. Folglich
ist kein Produkt-Singletonbild eine Einermenge.

Soll \(E\) unter diesem \(\Phi\) aus Markierungen mit
Singletonbildern bestehen, muss \(E=\varnothing\) gelten. Dann sind
alle Profile leer und die einzige Kontextklasse ist \(L\).
Sie enthält drei verschiedene Produktwerte. Die Voraussetzung
von Satz~\ref{cp:separation} ist also für dieses vorgegebene
\(\Phi\) falsch.

Das Hindernis betrifft die gewählte Abbildung: Die Grundhalbgruppe
ist natürlich zu sich selbst isomorph, und die Identität auf
\(\Pp L\) würde sämtliche Produktwerte korrekt markieren.
Für die allgemeine Existenz eines geeigneten Ersatzes fehlt
jedoch genau die folgende Aussage.

\begin{satz}[Äquivalenz der Normalisierung]\label{cp:normalization}
Für feste endliche Halbgruppen \(C,D\) und einen Isomorphismus
\(\Phi:\Pp C\to\Pp D\) sind äquivalent:
\begin{enumerate}
 \item \(C\cong D\).
 \item Es existieren \(\alpha\in\operatorname{Aut}(\Pp D)\)
 und Markierungen \(E\subseteq C^2\), sodass
 \(\Psi=\alpha\circ\Phi\) die markierten Einermengen auf
 Produkt-Einermengen abbildet und die zugehörigen Kontextkongruenzen
 auf beiden Seiten die Produktwerte trennen.
\end{enumerate}
\end{satz}

\begin{proof}
Aus der zweiten Aussage folgt die erste durch
Satz~\ref{cp:separation}. Umgekehrt induziert jeder
Grundisomorphismus \(f:C\to D\) den Potenzisomorphismus
\(\Pp f:A\mapsto f[A]\). Setze
\[
 \alpha=(\Pp f)\circ\Phi^{-1}.
\]
Dann ist \(\alpha\) ein Automorphismus von \(\Pp D\) und
\(\Psi=\Pp f\). Wähle \(E=C^2\). Alle Produktwerte werden als
Einermengen zugeordnet, und jeder ist bereits durch den leeren
Kontext von allen anderen Elementen getrennt.
\end{proof}

Die allgemeine Existenz einer solchen Normalisierung ist daher
logisch äquivalent zur allgemeinen endlichen Vermutung.
Der Profilsatz macht diese Voraussetzung präzise; er beweist
ihre allgemeine Erfüllbarkeit noch nicht.

\section{Ein Hindernis bei vier Produktwerten}

Die schlechte Einermengentreue ist nicht auf Linksnullhalbgruppen
beschränkt. Ein Beispiel zeigt die nächste Grenze selbst bei
verschwindenden Dreifachprodukten.

Sei \(W=\{0,u,v,w\}\), \(X=\{a,b,c\}\) und
\(S=W\mathbin{\dot\cup}X\). Definiere
\(\rho(a)=u,\rho(b)=v,\rho(c)=w\) sowie
\[
 xy=\rho(y)\quad(x,y\in X).
\]
Alle Produkte mit einem Faktor aus \(W\) seien Null.
Dann ist \(S\) assoziativ mit \(S^2=W\) und \(S^3=\{0\}\).

Sei \(\sigma\) die Permutation von \(\Pp X\), welche die Paare
\[
 \{a\}\leftrightarrow\{a,b\},\quad
 \{b\}\leftrightarrow\{a,c\},\quad
 \{c\}\leftrightarrow\{b,c\}
\]
vertauscht und \(X\) festlässt. Für \(A\ne\varnothing\),
\(A\subseteq X\), \(Z\subseteq W\) setze
\[
 \Phi(A\cup Z)=\sigma(A)\cup Z.
\]
Auf nichtleeren Teilmengen von \(W\) definieren wir
\(\Phi(\{0\})=\{0\}\) und
\[
 \Phi(Z_0\cup\rho[B])=Z_0\cup\rho[\sigma(B)]
 \quad
 (Z_0\subseteq\{0\},\ \varnothing\ne B\subseteq X).
\]
Beide Vorschriften betreffen disjunkte Bereiche und definieren
zusammen eine Involution.

Sind die beiden aktiven Teile \(A,B\) nichtleer, ist das Produkt
zweier dekorierter Mengen gleich \(\rho[B]\), ergänzt um \(0\)
genau dann, wenn mindestens eine ihrer \(W\)-Dekorationen
nichtleer ist. Diese Regel wird von \(\Phi\) erhalten.
Fehlt auf einer Seite der aktive Teil, ist das Produkt \(\{0\}\).
Also ist \(\Phi\) ein Potenzautomorphismus. Er erfüllt
\[
 \Phi(\{u\})=\{u,v\},\quad
 \Phi(\{v\})=\{u,w\},\quad
 \Phi(\{w\})=\{v,w\}.
\]
Von vier Produkt-Einermengen bleibt nur \(\{0\}\) einelementig.
Die drei für Folgerung~\ref{cp:kminusone} hinreichenden
Markierungen stehen unter diesem gegebenen Automorphismus nicht
zur Verfügung.

\section{Welcher Induktionsschritt noch fehlt}

Fasst \(\equiv_E\) verschiedene Produktwerte zusammen, so gilt
tatsächlich
\[
 \bigl|(C/{\equiv_E})^2\bigr|
 =|\pi_C(C^2)|<|C^2|.
\]
Die Induktionsgröße wird kleiner. Zugleich kennen wir aus
Satz~\ref{cp:profiles} den Quotientenisomorphismus bereits direkt.
Die zusätzliche Anwendung von \(\mathsf H(k-1)\) liefert keine
Zuordnung der verschiedenen Produktwerte innerhalb einer
zusammengefassten Klasse.

Für Grundklassen \(K,L\) fehlen gegebenenfalls die einzelnen
Multiplikationsabbildungen
\[
 \mu_{K,L}:K\times L\longrightarrow
 \{c\in C^2:\pi_C(c)=KL\},\qquad (a,b)\longmapsto ab.
\]
Der Quotient bestimmt die Zielklasse \(KL\), die Gewichte
bestimmen die Größen der Klassen. Das Beispiel aus
\eqref{cp:example-classes} zeigt, dass diese Daten die
Abbildungen \(\mu_{K,L}\) nicht bestimmen.

Ein vollständiger Induktionsschritt müsste aus dem ursprünglichen
Potenzisomorphismus zusätzliche, miteinander verträgliche Daten
zur Rekonstruktion dieser Abbildungen gewinnen oder die
Normalisierung aus Satz~\ref{cp:normalization} allgemein beweisen.
Beides ist hier nicht hergeleitet. Daher steht der vollständige
Beweis von \(\mathsf H(k)\) für beliebiges \(k\) weiterhin aus.

Bewiesen sind der gewichtete Kontextquotient, die Rekonstruktion
bei getrennten Produktwerten, die Folgerung mit \(k-1\) Markierungen
und der kurze Dreiprodukt-Nullbeweis. Die Gegenbeispiele bestimmen
genau, welche stärkeren Behauptungen aus diesen Sätzen nicht
folgen.

\section*{Einordnung und Quellen}

Die allgemeine endliche Frage wird in Tringalis Arbeit
\emph{On the Isomorphism Problem for Power Semigroups}
(2025), Einleitung, als offen beschrieben. Der frei zugängliche
Preprint vom 16.~August 2024 enthält diese Aussage ausdrücklich:
\href{https://arxiv.org/html/2402.11475v2}{arXiv:2402.11475v2}.
Die Verlagsfassung erschien in \emph{Recent Progress in Ring and
Factorization Theory}, Springer Proceedings in Mathematics and
Statistics~477, S.~429--437,
\href{https://doi.org/10.1007/978-3-031-75326-8_20}{DOI: 10.1007/978-3-031-75326-8\_20}.
Diese Einordnung ersetzt keinen mathematischen Beweis der hier
untersuchten Teilaussagen.

Die Lesefassung baut auf dem bereitgestellten Kontextprofil-Entwurf
und dem bereits im Skript bewiesenen Zweiproduktsatz auf.
Die Beispiele wurden zusätzlich vollständig an ihren endlichen
Tafeln geprüft, einschließlich der leeren Kontexte. Die Beweise
oben begründen die Aussagen unabhängig von diesen Kontrollen.
Eine literaturgeschichtliche Neuheit wird nicht beansprucht.

\end{document}
